\documentclass[12pt]{article}
\usepackage[margin=1in]{geometry}
\usepackage{setspace}
\usepackage{amsmath,amssymb,amsthm}
\usepackage{booktabs}
\usepackage{enumitem}
\usepackage[hidelinks]{hyperref}
\usepackage{xcolor}
\usepackage{titlesec}
\usepackage{longtable}
\usepackage{array}
\setstretch{1.15}
\setlist[itemize]{leftmargin=1.5em,itemsep=0.3em,topsep=0.3em}
\setlist[enumerate]{leftmargin=1.7em,itemsep=0.3em,topsep=0.3em}
\titleformat{\section}{\large\bfseries}{\thesection.}{0.5em}{}
\titleformat{\subsection}{\normalsize\bfseries}{\thesubsection.}{0.5em}{}

\title{Memo: How to Incorporate Three Affective-Polarization Papers into the Auditor-Change Draft}
\author{Prepared for revision planning}
\date{\today}

\begin{document}
\maketitle

\noindent\textbf{Papers assessed.}
\begin{itemize}
    \item Iyengar, Lelkes, Levendusky, Malhotra, and Westwood (2019), \emph{The Origins and Consequences of Affective Polarization in the United States}.
    \item Iyengar, Sood, and Lelkes (2012), \emph{Affect, Not Ideology: A Social Identity Perspective on Polarization}.
    \item Tyler and Iyengar (2023), \emph{Testing the Robustness of the ANES Feeling Thermometer Indicators of Affective Polarization}.
\end{itemize}

\section{Bottom line}

These three papers are highly relevant to the current draft, but they should be incorporated \emph{selectively and asymmetrically}. They are most useful for four purposes:
\begin{enumerate}
    \item sharpening the paper's conceptual distinction between \emph{ideological} and \emph{affective} polarization;
    \item replacing the draft's still-too-narrow emphasis on heterogeneous priors with a broader mechanism based on \emph{identity, distrust, and differential interpretation};
    \item making the measurement discussion more disciplined by explaining what feeling thermometers do and do not capture;
    \item motivating a feasible affective-polarization \emph{extension} rather than a full state-year horse race that the current data cannot support.
\end{enumerate}

They should \emph{not} be used to reframe the paper as primarily a measurement paper, and they should \emph{not} induce claims that the current draft already separately identifies ideological and affective polarization at the same geographic and temporal level. That would overpromise relative to the likely data pipeline.

\section{What each paper contributes}

\subsection{Iyengar et al. (2012): affective polarization is not the same as issue polarization}

This paper is the strongest direct support for the draft's core conceptual move away from issue-distance alone. Its central claim is exactly the one your paper needs: polarization in the mass public should not be understood only in terms of policy extremity; partisan affect and social distance are distinct, consequential, and in many settings more diagnostic than ideology alone.

For your draft, the main takeaways are:
\begin{itemize}
    \item partisan identity can produce negative evaluations of the out-group even when ideological disagreement is limited;
    \item social-distance outcomes and trait stereotyping are valid manifestations of polarization, not just issue divergence;
    \item exposure to campaigns and negative messages can intensify affective hostility.
\end{itemize}

This paper is therefore directly useful in the \textbf{Introduction}, \textbf{Related Literature}, and \textbf{Theory} sections. It supports language such as: polarization can change how investors interpret auditor disclosures even when investors are not far apart on specific economic-policy issues, because the disclosure is filtered through partisan identity and institutional distrust.

\subsection{Iyengar et al. (2019): social-identity mechanism and downstream consequences}

The 2019 Annual Review article provides the best broad synthesis for your mechanism. It defines affective polarization as dislike and distrust of the other party, roots it in partisan social identity, and distinguishes it from ideological polarization. It also surveys measurement approaches: feeling thermometers, trait ratings, social distance, implicit measures, and behavioral discrimination.

For your paper, this article is especially useful because it does three things the draft currently needs:
\begin{itemize}
    \item it gives a clean conceptual definition of affective polarization,
    \item it explains why affective polarization can rise even when issue disagreement is not the main driver,
    \item it shows that polarization spills over into nonpolitical settings, which makes it more plausible that it also spills over into the interpretation of auditor-originated disclosures.
\end{itemize}

This article should be the main reference when the draft says that polarization can operate through \emph{institutional distrust}, \emph{social identity}, and \emph{out-group animus}, not just through differences in numeric priors.

\subsection{Tyler and Iyengar (2023): a practical warning about measurement}

This paper is conceptually useful but methodologically even more important. It evaluates the ANES feeling-thermometer series and asks whether the standard affective-polarization measure is overstated by survey design. The key conclusion is nuanced: the increase in out-party animus is real, but mixed survey modes can inflate the time series, while selection bias and priming are less concerning than some critics suggest.

For your draft, this paper matters for two reasons:
\begin{itemize}
    \item it legitimizes use of feeling-thermometer-based affective measures as a serious benchmark rather than a casual proxy;
    \item it forces the paper to be careful about claiming that an ANES-based affective series is a perfectly clean annual measure.
\end{itemize}

So this paper should not be used merely as another citation for ``affective polarization exists.'' It should instead discipline the measurement section: if you use an ANES-based affective series, you should acknowledge that it is a national survey-based proxy, not a direct firm-level or state-year observation, and that mode shifts matter especially in the later period.

\section{How these papers should change the draft's substantive framing}

\subsection{The mechanism should no longer center on heterogeneous priors alone}

The current draft says that polarization increases heterogeneity in investors' priors, institutional trust, and perceived signal precision. That is much better than a pure Bayesian-priors story, but the writing still places too much emphasis on priors as the formal mechanism and too little on identity and distrust as the real economic story.

After incorporating these three papers, the draft should state the mechanism this way:

\begin{quote}
Political polarization increases heterogeneity in how investors interpret auditor-originated disclosures by intensifying partisan social identity, distrust of out-group-linked institutions, and disagreement over the credibility and precision of institution-dependent signals.
\end{quote}

That formulation is superior because:
\begin{itemize}
    \item it matches Iyengar et al. (2012) and Iyengar et al. (2019),
    \item it keeps the Bayesian model as a tractable microfoundation rather than the only plausible behavioral story,
    \item it better fits the audit setting, where trust in gatekeepers is central.
\end{itemize}

\subsection{The draft should explicitly distinguish ideological from affective polarization}

The paper should distinguish the two concepts clearly:
\begin{itemize}
    \item \textbf{Ideological polarization}: divergence in substantive issue positions or electoral/elite distributions.
    \item \textbf{Affective polarization}: out-party animus, distrust, social distance, and identity-based hostility.
\end{itemize}

But the distinction should be presented as a \emph{conceptual map}, not as a claim that the current baseline empirics separately identify both dimensions with equal precision. The right claim is:

\begin{quote}
The baseline empirical design measures the ideological-political environment directly. The affective-polarization literature clarifies the broader mechanism and motivates a supplementary national measure or validation exercise, rather than a fully symmetric state-year horse race.
\end{quote}

That wording keeps the paper honest and implementable.

\section{Concrete drafting changes}

\subsection{Abstract}

The abstract should avoid sounding as though the paper is purely about policy disagreement or Bayesian updating. A revised mechanism sentence should read along the following lines:

\begin{quote}
We argue that political polarization increases heterogeneity in how investors interpret auditor disclosures by amplifying partisan identity, institutional distrust, and differences in perceived signal credibility.
\end{quote}

That sentence should replace or dominate any wording that foregrounds heterogeneous priors alone.

\subsection{Introduction}

Add one paragraph after the first statement of the research question. The paragraph should do three things:
\begin{enumerate}
    \item define affective polarization as out-party animus and distrust;
    \item distinguish it from ideological polarization;
    \item explain why auditor disclosures are especially exposed to both dimensions because they are standardized but interpretation-heavy and depend on trust in regulated gatekeepers.
\end{enumerate}

A suitable draft paragraph is:

\begin{quote}
Recent work in political science distinguishes ideological polarization from affective polarization. The former concerns divergence in issue positions; the latter concerns partisan animus, social distance, and distrust of the opposing side. This distinction is important for auditor-originated disclosures. Investors may disagree not only because they attach different substantive implications to the same disclosure, but also because they differ in the extent to which they trust auditors, oversight institutions, and information associated with actors they perceive as politically or socially distant.
\end{quote}

\subsection{Related literature}

The literature review should add a dedicated subsection titled something like:

\begin{quote}
\textbf{Affective Polarization, Social Identity, and Information Processing}
\end{quote}

This subsection should:
\begin{itemize}
    \item cite Iyengar et al. (2012) for the proposition that affective polarization is distinct from issue-based polarization;
    \item cite Iyengar et al. (2019) for the social-identity mechanism and the variety of outcome domains it affects;
    \item cite Tyler and Iyengar (2023) for the robustness and limits of the ANES feeling-thermometer series.
\end{itemize}

This section should not read like a survey of political science for its own sake. It should be tied directly to the paper's setting: the interpretation of auditor changes is vulnerable to identity-based distrust because the disclosure is institution-dependent and often ambiguous.

\subsection{Theory section}

The theory section should explicitly say that the formal model captures a broader mechanism. The draft currently gestures in that direction, but it should be made more pointed. Add a paragraph such as:

\begin{quote}
The formal model operationalizes heterogeneity through differences in priors, but the broader mechanism is not limited to prior beliefs. In the affective-polarization literature, partisan identity and out-party distrust shape how individuals evaluate the credibility of information and the institutions that produce it. In the present setting, these forces can be interpreted as shifting priors, perceived signal precision, or trust in the gatekeeping function of auditors and regulators. The empirical predictions are similar across these interpretations: greater polarization yields greater dispersion in reactions to the same auditor-related disclosure.
\end{quote}

This directly integrates the three papers without forcing you to rebuild the theory.

\subsection{Measurement section}

This is where the current draft needs the most discipline.

The paper should add a short subsection titled:

\begin{quote}
\textbf{Ideological and Affective Polarization: Measurement Strategy}
\end{quote}

That subsection should explain:
\begin{enumerate}
    \item the baseline state-year measure is ideological/political rather than affective;
    \item affective polarization is conceptually central but harder to observe at the firm-event level;
    \item national affective-polarization measures based on ANES feeling thermometers are plausible supplementary measures, but they must be treated as national survey-based proxies and interpreted cautiously because of mode effects.
\end{enumerate}

The draft should \emph{not} say that Esteban--Ray and ANES thermometers are simply two alternative measures of the same object. They are not. One is closer to ideological/electoral clustering; the other is explicitly affective and survey-based.

\section{Feasible empirical use of these papers}

\subsection{What is feasible now}

A practical and defensible design is:
\begin{itemize}
    \item \textbf{Main specification}: use the existing ideological state-year measure around Item 4.01 events.
    \item \textbf{Supplementary affective test}: build a national ANES-based affective-polarization series using the in-party minus out-party feeling-thermometer differential, then interact it with a cross-sectional exposure variable and event ambiguity.
\end{itemize}

This structure aligns with Tyler and Iyengar (2023), because it treats the ANES series as a usable but not perfect benchmark. It also avoids pretending that you already have a high-quality state-year panel of affective polarization stretching cleanly across the full sample.

\subsection{What is not feasible to claim now}

The draft should avoid all of the following claims unless the data are actually built:
\begin{itemize}
    \item that affective polarization is measured at the same geographic level as the ideological measure;
    \item that the paper cleanly horse-races ideological versus affective polarization in the main specification;
    \item that the ANES series can be treated as a frictionless annual panel without survey-design caveats.
\end{itemize}

Tyler and Iyengar (2023) is especially useful here because it gives the paper a principled reason to be cautious rather than evasive.

\section{How to cite and deploy each paper in the manuscript}

\subsection{Iyengar, Sood, and Lelkes (2012)}

Use this paper in three places:
\begin{itemize}
    \item \textbf{Introduction}: to justify the claim that polarization should not be reduced to issue distance.
    \item \textbf{Literature Review}: to motivate social-distance and trait-based notions of polarization.
    \item \textbf{Mechanism discussion}: to support the claim that campaigns and partisan cues intensify out-party animus.
\end{itemize}

\subsection{Iyengar et al. (2019)}

Use this paper as the core synthesis reference in:
\begin{itemize}
    \item \textbf{Introduction}: for the conceptual distinction between ideological and affective polarization;
    \item \textbf{Theory}: for the social-identity account of why distrust and hostility matter;
    \item \textbf{Conclusion}: for the idea that affective polarization has spillovers outside formal politics, making capital-market interpretation effects plausible.
\end{itemize}

\subsection{Tyler and Iyengar (2023)}

Use this paper in:
\begin{itemize}
    \item \textbf{Measurement Strategy}: to justify thermometer-based measures while acknowledging design sensitivity;
    \item \textbf{Robustness Discussion}: to explain why any ANES-based affective measure should be interpreted as a supplementary national validation test, not as a perfect baseline.
\end{itemize}

\section{Recommended revision language}

\subsection*{Recommended sentence for the introduction}

\begin{quote}
The mechanism we study is not limited to ideological disagreement over policy implications. Affective polarization---the tendency of partisans to distrust and negatively evaluate the opposing side---can also shape how investors assess the credibility of auditors and the institutions that regulate them.
\end{quote}

\subsection*{Recommended sentence for the theory section}

\begin{quote}
Although the formal model represents heterogeneity through priors, the broader mechanism also encompasses identity-driven distrust and differences in perceived signal precision, both of which are central to the modern literature on affective polarization.
\end{quote}

\subsection*{Recommended sentence for the measurement section}

\begin{quote}
Our baseline empirical design measures the ideological-political environment surrounding the firm. We treat affective polarization as conceptually central but empirically harder to observe at the same level of aggregation, and therefore use survey-based feeling-thermometer measures, when employed, as supplementary national proxies rather than as full substitutes for the baseline state-year measure.
\end{quote}

\section{Overall recommendation}

These three papers should be incorporated. They materially improve the draft's conceptual precision and make the mechanism more persuasive. But they should be used to \emph{discipline} the paper, not to overcomplicate it.

The strongest revised version of the paper would:
\begin{enumerate}
    \item define affective polarization clearly and distinguish it from ideological polarization;
    \item broaden the mechanism from heterogeneous priors to identity, distrust, and perceived credibility;
    \item keep the current ideological measure as the empirical baseline;
    \item add, if feasible, a carefully caveated national affective-polarization extension using ANES feeling thermometers;
    \item avoid claiming a symmetric ideological-versus-affective horse race unless the data truly support it.
\end{enumerate}

That is the most credible way to integrate these three papers into the current auditor-change draft.

\end{document}
